\section{Variable-Granularity Object Storage}
\label{sec:layout}

A file read names a logical byte range, but its cost depends on the objects
containing that range, the bytes fetched, and the storage units verified.
Large objects can reduce request count for sequential access while increasing
the work needed to serve a partial read. \sysname{} therefore separates
client object formation from backend encoding and verification. The
connection between them is the finalized object, whose length becomes known
after assembly, applicable compression, and framing.

\subsection{From file ranges to objects}

The ordinary data path maps file offsets to logical chunks and resolves
the visible slices within each chunk. A chunk is an addressing range;
slices identify the written content that supplies its visible intervals.
The uploader partitions each slice's data using the configured client
block size $g$ and stores the resulting blocks under the slice identity
and block index. A 10 MiB slice with $g=4\MiB$, for example, can produce
raw blocks of 4, 4, and 2 MiB.

Compression, headers, and framing then determine each object's finalized
upload length $L_o$. Thus $g$ controls the partition of the slice, whereas
$L_o$ describes the bytes actually uploaded. Flush boundaries and partial
writes affect which slices and objects form, and a short tail need not
occupy a full block. When reading, the client uses the same geometry to
translate a visible slice interval into object references and offsets.
A new slice can replace one interval while untouched intervals continue
to reference older objects; the replacement need not rewrite the complete
old object.

The ordinary path defaults to 64 MiB chunks and $g=4\MiB$, with a chunk
size that is a nonzero multiple of $g$. This is a fixed configuration.
The proposed 1/4/8/16 MiB choices define an object-formation design space,
not an implemented controller that selects sizes from access behavior.
Because readers reconstruct block addresses from the configured geometry,
changing $g$ for existing data requires a compatible representation or
migration.

\subsection{Mapping objects to backend storage}

Once $L_o$ is known, a compatible RustFS research backend can use it to
inform a separate layout decision. Let $B$ denote the backend's
erasure-coding block size. The client may request $B$ through an
encoding-block-size hint $h(L_o)$, without changing $g$ or multipart part
sizes. The backend validates the request under its compatibility and
eligibility rules and persists the effective $B$. Omitting a hint retains
default selection. Multipart initiation uses the complete object's length,
and retries retain that object's decision.

The current interface supports fixed hints of 64 KiB, 256 KiB, 1 MiB, and
4 MiB. Its automatic hint policy omits the hint for $L_o\leq1\MiB$ and
requests $B=1\MiB$ otherwise. This rule acts on the formed object; it does
not select the client block size. The backend also defaults to $B=1\MiB$,
so automatic and default settings need not produce different layouts.

The physical verification/read quantum $q$ determines how much shard data
must be verified together.
For $d$ data shards and $p$ parity shards, let $s$ be one complete stripe's
logical shard size. A separately enabled backend policy can select a
physical verification/read quantum for eligible new framing subject to
\begin{equation}
q\mid s, \qquad qd\mid B.
\end{equation}
These are admission constraints for the new framing, preserving stripe
alignment while permitting smaller verification units. Both hint admission
and new $q$ framing have their own compatibility gates. Reads and healing
use persisted object geometry even after write settings change.

\subsection{From an object range to physical work}

Object size alone does not determine the bytes fetched for a file read.
Range GET can retrieve part of an object when its format permits independent
access. The ordinary store supports page-aligned range reads for eligible
small requests; compressed blocks or a whole-object fetch policy can require
larger transfers. A packed frame provides a decoding and integrity boundary
inside a snapshot object, distinct from both the client block size $g$ and
the backend quantum $q$. Caching and concurrent-read coalescing determine
which of these requests reach the backend.

The backend maps the resulting object range into shard intervals.
For a shard interval $[a,a+\ell)$, with $a\geq0$, $\ell>0$, and $q>0$,
the number of verification/read quanta intersected is
\begin{equation}
n_q(a,\ell)=
\left\lceil\frac{a+\ell}{q}\right\rceil-
\left\lfloor\frac{a}{q}\right\rfloor.
\end{equation}
Ignoring tails and batching, the covered data is approximately $n_q q$
plus checksums, summed over shards actually read. The application request
cannot be substituted directly for $\ell$: the client range, compression,
and coding determine that interval. Smaller $q$ can reduce over-read while
increasing checksum overhead. For fixed $(d,p)$, the nominal redundancy
remains $(d+p)/d$; a large $B$ does not imply a fully padded tail stripe.

\subsection{Costs and scope}

The tradeoff in $g$ is visible even before backend encoding.
Table~\ref{tab:object-granularity} considers one uncompressed 64 MiB slice
starting at file offset zero and a read of $[10,12)\MiB$. Increasing $g$
reduces the number of objects in the slice, but increases bytes fetched
for this read if complete objects must be obtained. These are arithmetic
examples, not measurements or a policy recommendation.

\begin{table}[t]
\centering
\begin{tabular}{rrr}
\toprule
$g$ & Object count & \shortstack{Whole-object\\bytes fetched}\\
\midrule
1 MiB & 64 & 2 MiB\\
4 MiB & 16 & 4 MiB\\
8 MiB & 8 & 8 MiB\\
16 MiB & 4 & 16 MiB\\
\bottomrule
\end{tabular}
\caption{Idealized object formation and a 2 MiB partial read. The last
column assumes complete fetches of the intersecting objects, with no cache
hits or compression and ignoring framing overhead; it does not describe
every range-read path.}
\label{tab:object-granularity}
\end{table}

These layers require separate cost accounting. Network read amplification
compares object response bytes with application bytes served; upload
amplification compares uploaded bytes with application bytes written;
physical read amplification compares storage-device bytes with application
bytes served. Request counts, compression work, and cache occupancy capture
additional costs. Compaction, retries, and accumulated short slices can
increase upload work, so smaller objects do not guarantee lower write
amplification. Similarly, whole-object fetching can prevent a smaller $q$
from reducing the work generated by a partial application read.

The mechanism applies to an ordinary flat volume without delta-based
workspaces or Packed Metadata. Either of those mechanisms can also retain
fixed client geometry and default backend layouts. Their composition
changes the referenced objects and access patterns, while snapshot
producers and workspace upper writers retain their own object formats and
upload contracts. Section~\ref{sec:implementation} describes these path
boundaries and the ordinary client's coupling of enabled hints with
whole-object fetch and cache policy. That coupling must be separated from
backend geometry when attributing a measured effect to $B$ or $q$.
